\section{Orthogonal Stress-Matched Control Panel (Pre-Registered Extension)}
\label{sec:stresspanel}

A surviving alternative explanation for the primary result is broad
stress-adaptation rather than radiation-specific selection. To close this
gap, a stress-matched, radiation-sensitive control panel was locked by
amendment before any control data were downloaded or inspected.

\subsection{Panel definition (locked)}
Four species, all stress-tolerant and radiation-sensitive by prior
literature: \emph{Pyrococcus furiosus} (hyperthermophile; reference
assembly GCF\_008245085.1), \emph{Sulfolobus solfataricus}
(thermoacidophile; GCF\_003852155.1; NCBI taxonomic reclassification
\emph{Saccharolobus solfataricus} noted in the provenance ledger),
\emph{Lactiplantibacillus plantarum} (acid- and osmo-tolerant;
GCF\_009913655.1), and \emph{Saccharomyces cerevisiae} (osmo- and
ethanol-tolerant; GCF\_000146045.2). The amendment further locked that if
any panel member proves radiation-resistant, it remains in the panel and
the fact is reported; no swapping is permitted.

\subsection{Identical pipeline}
The control proteomes were appended to the tagged proteome set (168,186
proteins across the expanded panel) and re-clustered with the identical
locked orthogroup-calling pipeline: cascaded sensitive MMseqs2 clustering,
minimum sequence identity 0.3, minimum coverage 0.5, cov-mode 1. All formal
tests (one-sided Fisher presence/absence, one-sided Mann--Whitney U on copy
number, BH FDR 0.05, 1000 label permutations, $\geq$2-of-3 independent
origins rule) are re-run unchanged on the expanded clustering.

\subsection{Specificity layer (locked threshold)}
Each of the twelve gate-passing candidate orthogroups is re-tested for
radiation specificity: the radiation-specificity claim survives only for
candidates absent from at least three of the four stress-matched controls.
Candidates failing this threshold are re-labeled broad-stress-adaptation,
and the paper reports them as such.
